%
% Portada.
%

% Nota: Será más cómodo emplear el comando \maketitle que genera una portada de forma automática, pero 
% no incluye toda la información que es necesario incluir en la memoria de un proyecto de fin de carrera
% de la Facultad de Informática de A Coruña.
%

\graphicspath{ {Graphics/} }

\begin{titlepage}

	\begin{center}

		% Logotipo de la universidad.
	\includegraphics[width = 8cm] {./Graphics/logocei.eps}
		\vspace{1cm}

		% Nombre de la facultad, de la universidad y del departamento en que se realiza el PFC.
		{\Large{\textbf{UNIVERSITY OF INFORMATICS OF MADRID - UPM}}}
		\\
		{\it \large{\textbf{ARTIFICIAL INTELLIGENCE DEPARMENT}}}
		\vspace{1cm}

		% Indicamos el nombre de la titulación oficial que hemos cursado con tanto esfuerzo.
		{\large INTELLIGENT SERCHING BASED ON METHAHEURISTICS\\MASTER IN ARTIFICIAL INTELLIGENCE}
		\vspace{1cm}

		% Título
		\textbf{\Large Solving the magic square problem with the NSGA II genetic algorithm}
		\vspace{6cm}
	\end{center}

	\begin{flushright}
		\begin{tabular}{ll}
			% Nombre del alumno.
			\large{\textbf{Alumni:}}	&
			\large{Petra Franjic, Gustavo Puig} \\

			% Nombre del director/tutor del proyecto.
			\large{\textbf{Professor:}}	&
			\large{D. Alfonso Mateos Caballero} \\

			% Fecha.
			\large{\textbf{Date:}}	&
			\large{2th November 2015} \\
		\end{tabular}
	\end{flushright}

\end{titlepage}
